% Chapter III, §35 and Bibliographical Notes
% Chandrasekhar, "Hydrodynamic and Hydromagnetic Stability"
% Transcribed from scanned PDF

%----------------------------------------------------------------------
\section{Experiments on the onset of thermal instability in rotating fluids}
\label{sec:35}
%----------------------------------------------------------------------

We have seen that rotation influences the onset of thermal instability
in many ways.  The principal effects are these.

\begin{enumerate}
\item[(1)] Rotation inhibits the onset of instability; the extent of the
  inhibition depends on the Taylor number $\mathscr{T}$ ($= 4\Omega^2 d^4/\nu^2$) and the
  Prandtl number $p$ ($= \nu/\kappa$); and as $\mathscr{T} \to \infty$, the Rayleigh number at which
  instability sets in, and the wave number $a$ (in units of $1/d$) with which it
  is associated, have the asymptotic behaviours
  \begin{equation}
    R \to \text{constant}\; \mathscr{T}^{\!\nicefrac{2}{3}}
    \quad \text{and} \quad
    a \to \text{constant}\; \mathscr{T}^{\!\nicefrac{1}{6}}.
    \label{eq:3-308}
  \end{equation}

\item[(2)] The onset of instability will be as stationary convection so long
  as $p$ exceeds a certain critical value $p^*$.  The precise value of $p^*$ depends
  on the nature of the bounding surfaces; but it can be inferred from the
  value ($= 0.6766$) determined for the case of two free boundaries that
  $p^* \sim 1$.  When $p > p^*$, the $(R_c, \mathscr{T})$-relation is independent of $p$.

\item[(3)] The onset of instability will be as overstable oscillations if $p < p^*$
  and the Taylor number exceeds a certain value $\mathscr{T}^{(0)}$ depending on $p$.  For
  $\mathscr{T} \leq \mathscr{T}^{(0)}$ the onset of instability will be as stationary convection.  For a
  given $p$ ($< p^*$) the asymptotic behaviours of $R$ and $a$ are again as in (1)
  above; but the constants of proportionality now depend on $p$.  Also, the
  gyration frequency $p$ of the overstable oscillations is essentially deter\-mined
  by $\Omega$ and the Taylor number; for $0 < p < p^*$,
  \begin{equation}
    \nu/\Omega \to \text{constant}\; \mathscr{T}^{-\nicefrac{1}{6}}
    \quad \text{as} \quad \mathscr{T} \to \infty,
    \label{eq:3-309}
  \end{equation}
  where the constant depends on $p$.
\end{enumerate}

From the foregoing summary of the principal effects to be expected,
it is clear that in selecting the liquids for the experiments it would be
preferable to choose two liquids, one with a Prandtl number substantially
larger than unity and the other with a Prandtl number substantially less
than unity.  Two such liquids are readily available: water with $p \sim 7{\cdot}5$
and mercury with $p = 0{\cdot}025$.  Experiments with these two liquids have
been carried out by Nakagawa and Frenzen, Fultz and Nakagawa, and
Goroff.  An account of these experiments will now be given.

%----------------------------------------------------------------------
\subsection{Experiments with water}
\label{sec:35a}
%----------------------------------------------------------------------

In these experiments, water was contained in Pyrex glass cylinders
of 12 inches outside diameter.  To the bottoms of the cylinders were
cemented electrically conducting Pyrex heating plates; they served as
the heating elements to which measured amounts of electric power could
be supplied.  The cylinders containing the fluid and the necessary
accessories were placed on a supporting steel ring mounted on a vertical
shaft.  The steel ring had provisions for levelling and for other adjust\-ments.
The fluid was set in rotation by driving the shaft with suitable
belts and motors.  In the experiments, layers of water of depths varying
from 3 to 17~cm and speeds of rotation up to 60~rev/min were used.  Taylor
numbers up to $10^8$ were reached in this manner.

The difference in the temperatures at two levels, one near the top and
one near the bottom surface of the fluid, was measured by copper--constantan
wire thermocouples.  By using several junctions in series as
a thermopile, it was possible to measure the difference in the average
temperatures at the two levels with an accuracy of $\pm 0{\cdot}01$~$^\circ$C; this
accuracy in the temperature measurements was found to be necessary
and sufficient.

Observations with the foregoing experimental arrangement were of
two kinds: visual and photographic observations made with the aid of a
rotoscope which enabled the direct detection of the fluid motions
relative to the rotating frame; and quantitative observations based on
the records of the difference of temperature between the top and bottom
surfaces made continuously during each experiment after varying
amounts of electric power had been turned on to the heating elements.

By visual observations with the rotoscope, it was possible to observe
the onset of cellular convection when the heating rate exceeded a certain
amount.  The motions could be made visible by scattering a small
amount of aluminium powder on the top surface.  The movements of the
aluminium particles could be photographed by a time-exposure streak
technique.  An example of such a photograph is shown in Fig.~30.  The
internal motions in the cells could also be followed with black ink and

\figurePlaceholder{30}{Convection cells which appear in water when
in rotation and heated from below.  The data to which
the illustration refers are: depth 18~cm, difference in
temperature $0{\cdot}7^\circ$; rate of rotation $5{\cdot}0$~rev/min; Taylor
number $1{\cdot}2 \times 10^8$.}

\figurePlaceholder{31}{Side view of ink rising from the bottom of a
cylinder of water in the ascending cores of cells.  The
data to which the illustration refers are: depth 18~cm,
difference in temperature $0{\cdot}5^\circ$; rate of rotation $10{\cdot}9$~rev/min;
Taylor number $5{\cdot}5 \times 10^9$.}

\noindent photographed.  Fig.~31 is an example of a photograph showing such
internal motions.

For a quantitative determination of the critical Rayleigh number for
the onset of instability, a careful study of the temperature records is
necessary.  In Fig.~32 examples of temperature records which were
obtained for varying heating rates are shown.  Fig.~32\,$a$ corresponds to

\figurePlaceholder{32}{Time records of the adverse temperature gradient for water for three
different rates of heating ($d = 3$~cm, $\Omega = 10$~rev/min) (from \emph{Proc.\ Roy.\ Soc.\
(London)} A, \textbf{231}, 519 (1955)).}

\noindent a case when the bottom plate was heated at a sufficiently low rate that
the steady difference in temperature attained was insufficient to induce
instability; Fig.~32\,$b$ is a typical record when the heating rate sufficed to
surpass by some margin the temperature gradient necessary for the
onset of instability; Fig.~32\,$c$ is an example when the heating rate was so
much in excess of what is necessary to cause instability that the conditions
were approaching turbulence.  The difference between the records, when
we are in the conductive r\'egime and when we have passed it, is sufficiently
striking that by examining them for various heating rates, the tempera\-ture
gradient at which conditions are just past marginal could be deter\-mined
with some precision.  Simultaneous visual observations confirmed
the criteria that were used.  In this manner, Nakagawa and Frenzen

\figurePlaceholder{33}{Comparison of the theoretical relation with the experimental results
on the onset of thermal instability in water in varying speeds of rotation: the
critical Rayleigh number $R_c$ is plotted against the Taylor number $\mathscr{T}$.  Experi\-ments
with different depths of water are distinguished: $\boxtimes$ 18~cm; $\diamondsuit$ 14~cm;
$\odot$ 10~cm; $\bigotimes$ 6~cm; $\times$ 17.3~cm; $\triangle$ 13.2~cm; $\blacksquare$ 9.3~cm; $\circ$ 6.3~cm; $\boxminus$ 3~cm;
$\bigcirc$ 2~cm.}

\noindent determined the critical Rayleigh number as a function of the Taylor
number.  Their results are shown in Fig.~33.  The theoretical $(R_c, \mathscr{T})$-relation
for one rigid and one free surface derived in \S\,27\,$(c)$ is shown in
the same plot; it will be observed that the theoretically predicted relation
is fully confirmed.

Estimates of the cell dimensions based on photographic observations
are also in accord with the theoretical expectations.

%----------------------------------------------------------------------
\subsection{Experiments with mercury}
\label{sec:35b}
%----------------------------------------------------------------------

The experiments of Fultz and Nakagawa with mercury were carried
out with essentially the same experimental arrangement as described
in \S\,(a) above.  But a number of precautions peculiar to mercury had to
be observed.  For example, Fultz and Nakagawa record that considerable
trouble was experienced with the oxidation of the mercury surface.  It
was necessary on this latter account to replace the air above the heated
surface by nitrogen.  This was accomplished by circulating nitrogen;
this also served the purpose of keeping the temperature of the top surface
constant.  Even with these precautions, it was found that some time after
the heating current had been turned on, a thin contaminant film was
formed on the surface.  The film was rigid enough to prevent motions at
the top surface.  The layer of mercury was thus effectively confined
between two rigid surfaces.

The temperature gradients at which instability occurred were again
determined by an examination of the temperature records.  This was
supplemented by making plots of the temperature gradients against the
heating rates and using the Schmidt--Milverton principle (Chapter~II,

\figurePlaceholder{34}{Time records of the adverse temperature gradient for mercury for three
different rates of heating ($d = 6$~cm, $\Omega = 15$~rev/min) (from \emph{Proc.\ Roy.\ Soc.\
(London)} A, \textbf{231}, 520 (1955)).}

\noindent \S\,18\,(b)).  Layers of mercury of depth up to 8~cm and speeds of rotation
up to 30~rev/min were used.  Taylor numbers in the range $10^6$--$10^{12}$ were
thus accessible to the experiments.

In Fig.~34 a series of temperature records similar to those in Fig.~32
for water are shown.  The difference in the two sets of temperature
records, after instability has set in, is very striking; in the experiments
with mercury, the temperature records exhibit very pronounced periodic
fluctuations.  Fig.~35, exhibiting records which show these periodic
fluctuations with uniformity over durations as long as half an hour and
more, leaves very little doubt that we are witnessing here the overstable
oscillations predicted by the theory.  Indeed, by counting the number

\figurePlaceholder{35}{Further time records of the adverse temperature gradient for mercury.
The conditions to which they refer are given below:
\begin{center}
\begin{tabular}{lcc}
  & (a) & (b) \\
depth (cm) & 6 & 6 \\
$\Omega$ (sec$^{-1}$) & $1{\cdot}11$ & $3{\cdot}08$ \\
temperature difference (degrees) & $1{\cdot}54$ & $5{\cdot}05$ \\
Taylor number ($\mathscr{T}$) & $5 \times 10^9$ & $3{\cdot}9 \times 10^{10}$ \\
Rayleigh number ($R$) & $2{\cdot}64 \times 10^4$ & $8{\cdot}94 \times 10^4$ \\
period of oscillation (sec) & $19$--$20$ & $9$--$10$
\end{tabular}
\end{center}}

\noindent of these fluctuations in known intervals of time, we can estimate the
basic periods of the oscillations.  The results of such estimates are shown
in Fig.~36 together with the theoretically expected relation.  It will be
seen that the agreement is satisfactory.

\figurePlaceholder{36}{A comparison of the observed periods of oscillation at marginal stability
with the theoretical periods.  The ordinate $\tau$ gives the period in units of $2\varpi$.  The curves
$aa$ and $bb$ are the theoretical relations for $p = 0{\cdot}025$ and for the case of two bounding
surfaces rigid ($aa$) and one bounding surface rigid and the other free ($bb$).}

\figurePlaceholder{37}{A summary of the theoretical and experimental results on the critical Rayleigh
numbers for the onset of instability in mercury.  The curves $aa$, $bb$, and $cc$ are the
theoretical $(R_c, \mathscr{T})$-relations for the three cases (i) both bounding surfaces rigid, (ii) one
bounding surface rigid and the other free, and (iii) both bounding surfaces free for
$p = 0{\cdot}025$.  The solid circles are the experimentally determined points for mercury
($p = 0{\cdot}025$).}

In Fig.~37 the results on the critical Rayleigh numbers are shown
together with the predicted relations for the three sets of boundary
conditions.  It is noteworthy that the agreement is best with the
theoretical relation for two rigid boundaries: as we have already
remarked, the experimental conditions do correspond to this case.

The experiments of Nakagawa confirm that the onset of thermal

\figurePlaceholder{38a}{The variation of the Nusselt number for increasing Rayleigh numbers
for the case when the Taylor number is $1{\cdot}12 \times 10^9$.  The sharp break at the
Rayleigh number $3{\cdot}42 \times 10^4$ occurs while the system is already overstable.
Notice that the Nusselt number has not increased substantially beyond 1 even
though overstability set in at $R \sim 2{\cdot}5 \times 10^4$; this indicates the relative in\-efficiency
of heat-transfer by overstable oscillations.}

\figurePlaceholder{38b}{The dependence of the Rayleigh number on
the Taylor number at which the second break occurs.
The full-line curve is the theoretical relation for the
onset of stationary convection ignoring overstability.}

\noindent instability in a horizontal layer of mercury in rotation occurs as over\-stable
oscillations at the predicted Rayleigh numbers $R_c^{(o)}$ and with the
predicted characteristic frequencies.  When the Rayleigh number
exceeds $R_c^{(o)}$, the overstable oscillations will become of finite amplitude
and non-linear effects, ignored in the linear stability theory, will become
effective.  At the same time, the mode with the real exponential time
dependence will continue to be damped (on the assumption that the
presence of the overstable oscillations does not seriously alter the static
initial state to which the entire theory refers).  When $R$ becomes equal
to $R_c^{(s)}$ ($> R_c^{(o)}$), this latter mode will begin to manifest itself.  Since $R_c^{(s)}$
is about twenty-six times $R_c^{(o)}$ (for $\mathscr{T} \to \infty$, when both bounding surfaces
are rigid, and $p = 0{\cdot}025$), it is clear that the onset of stationary convec\-tion
will occur superposed on overstable oscillations of quite finite
amplitudes.  However, on general grounds, one may expect that the
overstable oscillations are not very efficient in transporting heat; and
that on this account one may, nevertheless, be able to detect the onset
of stationary convection in a Schmidt--Milverton plot.  Some experi\-ments
of Dropkin and Globe seemed to indicate that a second break\footnote{The first break was so weak in these experiments that one could not be certain of
it in most cases.}
occur in a Schmidt--Milverton plot after a first break corresponding to
the onset of overstability.  But these experiments were not carried out
with sufficient precautions to decide whether or not the second break
occurs at the predicted Rayleigh number $R_c^{(s)}$.  More recent experiments
by Goroff carried out with the necessary care confirms that the second
break does occur at the predicted place.  In Fig.~38\,$a$ the results of his
experiments for $\mathscr{T} = 1{\cdot}12 \times 10^9$ are shown; in this plot of the Nusselt
number against the Rayleigh number, a break clearly occurs at
$R = 3{\cdot}45 \times 10^4$; and this is long after overstable oscillations had become
a permanent feature of the temperature records.  The value
$R = 3{\cdot}45 \times 10^4$ should be compared with the theoretical value,
$R_c^{(s)} = 3{\cdot}50 \times 10^4$ which the results of Table~VIII predict for the onset
of stationary convection for $\mathscr{T} = 1{\cdot}12 \times 10^9$.  In Fig.~38\,$b$ the values of
$R_c^{(s)}$ obtained for three different values of $\mathscr{T}$ (including the one to which
Fig.~38\,$a$ refers) are shown together with the theoretically expected
relation.  It would appear from these experiments that there can be
little doubt that the onset of stationary convection occurs at the pre\-dicted
Rayleigh number.  They further establish the relative inefficiency
of overstable oscillations for the transport of heat.

%======================================================================
\section*{BIBLIOGRAPHICAL NOTES}
\label{sec:bib-notes-III}
%======================================================================

For a general account of the subject of this chapter see:

\begin{enumerate}
\item[1.] S.\ \textsc{Chandrasekhar}, `Thermal convection', \emph{Daedalus}, \textbf{86}, 323--39 (1957).
\end{enumerate}

\noindent \S\,22.  The fundamental theorem proved in this section is due to:
\begin{enumerate}
\item[2.] G.\ I.\ \textsc{Taylor}, `Experiments with rotating fluids', \emph{Proc.\ Roy.\ Soc.\ (London)}
  A, \textbf{100}, 114--21 (1921).
\end{enumerate}

\noindent The same theorem, though not in as explicit a form, was independently stated by:
\begin{enumerate}
\item[3.] J.\ \textsc{Proudman}, `On the motion of solids in a liquid possessing vorticity',
  \emph{Proc.\ Roy.\ Soc.\ (London)} A, \textbf{92}, 408--24 (1916).
\end{enumerate}

In Taylor's paper (reference~2), experiments are described which demonstrate
very effectively the implications of the theorem for the character of the fluid
motions which prevail under the circumstances.  Taylor's experiments have been
repeated by:
\begin{enumerate}
\item[4.] D.\ \textsc{Fultz}, `A survey of certain thermally and mechanically driven systems
  of meteorological interest', \emph{Proceedings of the First Symposium on the Use
  of Models in Geophysical Fluid Dynamics}, 27--63, edited by Robert B.\ Long,
  Johns Hopkins University, Baltimore, 1953.
\end{enumerate}

\noindent The following references to related papers by Taylor may be noted:
\begin{enumerate}
\item[5.] G.\ I.\ \textsc{Taylor}, `Motion of solids in fluids when the flow is not irrotational',
  \emph{Proc.\ Roy.\ Soc.\ (London)}, A, \textbf{93}, 99--113 (1917).
\item[6.] \rule{2em}{0.4pt}, `The motion of a sphere in a rotating liquid', ibid., \textbf{102}, 180--9
  (1922).
\item[7.] \rule{2em}{0.4pt}, `Experiments on the motion of solid bodies in rotating fluids', ibid.,
  \textbf{104}, 213--18 (1923).
\end{enumerate}

The dynamics of rotating fluids have many fascinating aspects of which the
Taylor--Proudman theorem is only one.  For a general account of these other
aspects see:
\begin{enumerate}
\item[8.] T.\ E.\ \textsc{Sterne}, `Rotating fluids', \emph{Surveys in Mechanics}, 139--61, edited by
  G.\ K.\ Batchelor and R.\ M.\ Davies, Cambridge Monographs on Mechanics
  and Applied Mathematics, Cambridge, England, 1956.
\end{enumerate}

\noindent See also:
\begin{enumerate}
\item[9.] G.\ W.\ \textsc{Morgan}, `A study of motions in a rotating liquid', \emph{Proc.\ Roy.\ Soc.\
  (London)} A, \textbf{206}, 108--30 (1951).
\item[10.] H.\ \textsc{G\"ortler}, `\"Uber eine Schwingungserscheinung in Fl\"ussigkeiten mit
  stabiler Dichteschichtung', \emph{Z.\ ang.\ Math.\ Mech.}\ \textbf{23}, 65--71 (1943).
\item[11.] \rule{2em}{0.4pt}, `Einige Bemerkungen \"uber Str\"omungen in rotierenden Fl\"ussigkeiten',
  ibid.\ \textbf{24}, 210--14 (1944).
\end{enumerate}

\noindent \S\S\,24--27.  The analyses in these sections are derived from:
\begin{enumerate}
\item[12.] S.\ \textsc{Chandrasekhar}, `The instability of a layer of fluid heated below and
  subject to Coriolis forces', \emph{Proc.\ Roy.\ Soc.\ (London)} A, \textbf{217}, 306--27 (1953).
\item[13.] \rule{2em}{0.4pt}\ and \textsc{Donna} D.\ \textsc{Elbert}, `The instability of a layer of fluid heated
  below and subject to Coriolis forces.\ II', ibid., \textbf{231}, 198--210 (1955).
\end{enumerate}

\noindent The method of solution described in \S\,27 (for cases $(b)$ and $(c)$) is, however, different
from that in reference~12; and the results included in Tables~VIII and IX (due
to Miss Donna Elbert) are based on the new formulae.

\noindent \S\,28.  The discussion of the cell patterns in this section follows:
\begin{enumerate}
\item[14.] G.\ \textsc{Veronis}, `Cellular convection with finite amplitude in a rotating fluid',
  \emph{J.\ Fluid Mech.}, \textbf{5}, 401--35 (1959).
\end{enumerate}

\noindent \S\,29.  The principal results of this section are contained in references~12 and~13.
But the arguments have been rearranged and made more explicit.  The particular
form of the discussion of the roots given in \S\,(a) is due to Dr.\ John Sykes.

\noindent \S\S\,30--32.  The analyses in these sections are again derived from references~12
and~13.

\noindent \S\,33.  See:
\begin{enumerate}
\item[15.] S.\ \textsc{Chandrasekhar}, `The thermodynamics of thermal instability in
  liquids', \emph{Max Planck Festschrift 1958}, 103--14, Veb Deutscher Verlag der
  Wissenschaften, Berlin, 1958.
\end{enumerate}

\noindent \S\,35.  The references for the experimental work described in this section are:
\begin{enumerate}
\item[16.] D.\ \textsc{Fultz}, Y.\ \textsc{Nakagawa}, and P.\ \textsc{Frenzen}, `An instance in thermal
  convection of Eddington's ``overstability''\,', \emph{Physical Rev.}, \textbf{94}, 1471--2
  (1954).
\item[17.] \rule{2em}{0.4pt}, `Experiments on overstable thermal convection in mercury',
  \emph{Proc.\ Roy.\ Soc.\ (London)} A, \textbf{231}, 211--25 (1955).
\item[18.] Y.\ \textsc{Nakagawa} and P.\ \textsc{Frenzen}, `A theoretical and experimental study of
  cellular convection in rotating fluids', \emph{Tellus}, \textbf{7}, 1--21 (1955).
\item[19.] D.\ \textsc{Dropkin} and S.\ \textsc{Globe}, `Effect of spin on natural convection in
  mercury heated from below', \emph{J.\ Appl.\ Phys.}, \textbf{30}, 84--89 (1959).
\item[20.] I.\ R.\ \textsc{Goroff}, `An experiment on heat transfer by overstable and ordinary
  convection', \emph{Proc.\ Roy.\ Soc.\ (London)} A, \textbf{254}, 537--41 (1960).
\end{enumerate}
